\documentclass[11pt,a4paper]{article}
\usepackage[italian]{babel}
\usepackage[margin=2.5cm]{geometry}
\usepackage{parskip}
\usepackage{amsmath, amssymb, amsthm}
\usepackage{booktabs}
\usepackage{array}
\usepackage{xcolor}
\usepackage{needspace}
\usepackage{enumitem}
\usepackage{listings}
\usepackage{tikz}
\usepackage[most]{tcolorbox}
\usepackage[hidelinks]{hyperref}
\usetikzlibrary{decorations.pathreplacing, shapes.geometric, arrows.meta, positioning}

% Niente vedove/orfane: i paragrafi non lasciano righe isolate a fine/inizio pagina
\widowpenalty=10000
\clubpenalty=10000

% Elenchi più compatti
\setlist{itemsep=2pt, parsep=0pt, topsep=4pt, partopsep=0pt}

% --- Riquadri con bordo nero, senza numero (si spezzano tra due pagine) ---
% Uso: \begin{definizione} ... \end{definizione}
%      \begin{teorema}[Nome] ... \end{teorema}  (il nome è facoltativo)
\tcbset{nota/.style={enhanced jigsaw, breakable, colback=white,
  colframe=black, boxrule=0.5pt, sharp corners}}
\NewTColorBox{definizione}{}{nota,
  before upper={\textbf{Definizione.}\ }}
\NewTColorBox{teorema}{o}{nota,
  before upper={\textbf{Teorema\IfValueT{#1}{ (#1)}.}\ }}
\NewTColorBox{proposizione}{o}{nota,
  before upper={\textbf{Proposizione\IfValueT{#1}{ (#1)}.}\ }}
\NewTColorBox{proprieta}{o}{nota,
  before upper={\textbf{Proprietà\IfValueT{#1}{ (#1)}.}\ }}
\NewTColorBox{assioma}{o}{nota, boxrule=1pt,
  before upper={\textbf{Assioma\IfValueT{#1}{ (#1)}.}\ }}

% --- Ambienti semplici (senza riquadro) ---
% Il titolo sta in cima al blocco, anche se il contenuto inizia con un riquadro o
% con codice; e non resta mai da solo in fondo a una pagina.
% Uso: \begin{esempio}[Titolo facoltativo] ... \end{esempio}
\newcommand{\newrunin}[2]{%
  \NewDocumentEnvironment{#1}{o}{%
    \par\Needspace{4\baselineskip}\addvspace{\medskipamount}\noindent
    \textbf{#2}\IfValueT{##1}{ (##1)}\textbf{.}\ \ignorespaces}%
  {\par\addvspace{\medskipamount}}}
\newrunin{esempio}{Esempio}
\newrunin{esercizio}{Esercizio}
\newrunin{osservazione}{Osservazione}

% V e F nelle tavole di verità
\newcommand{\V}{\textsf{V}}
\newcommand{\F}{\textsf{F}}

% --- Codice C ---
% Uso: \begin{codice} ... \end{codice}      codice C numerato
%      \begin{comando} ... \end{comando}    comandi da terminale
%      \begin{risultato} ... \end{risultato} output di un programma
%      \lstinline|printf("ciao");|           codice dentro al testo
\lstdefinestyle{c}{
  language=C,
  basicstyle=\ttfamily\small,
  keywordstyle=\color{blue}\bfseries,
  commentstyle=\color{gray}\itshape,
  stringstyle=\color{red!70!black},
  numbers=left,
  numberstyle=\tiny\color{gray},
  numbersep=8pt,
  columns=flexible,
  keepspaces=true,
  breaklines=true,
  showstringspaces=false,
  tabsize=4,
  morekeywords={include, define},
  % lettere accentate nei commenti
  literate={à}{{\`a}}1 {è}{{\`e}}1 {é}{{\'e}}1 {ì}{{\`i}}1
           {ò}{{\`o}}1 {ù}{{\`u}}1 {È}{{\`E}}1 {À}{{\`A}}1
}
\lstset{style=c}
\newtcblisting{codice}{listing only, nota, left=6mm,
  listing options={style=c}}
\newtcblisting{comando}{listing only, nota,
  listing options={basicstyle=\ttfamily\small, numbers=none, language={},
    columns=flexible, keepspaces=true, breaklines=true}}
\newtcblisting{risultato}{listing only, enhanced jigsaw, breakable,
  colback=black!5, colframe=black, boxrule=0.5pt, sharp corners,
  listing options={basicstyle=\ttfamily\small, numbers=none, language={},
    columns=flexible, keepspaces=true, breaklines=true}}

% --- Diagrammi di flusso (simboli standard) ---
\tikzset{
  terminale/.style = {ellipse, draw, minimum width=2.5cm, minimum height=0.9cm},
  io/.style        = {trapezium, trapezium left angle=70, trapezium right angle=110,
                      draw, minimum width=2.5cm, minimum height=0.9cm},
  processo/.style  = {rectangle, draw, minimum width=2.5cm, minimum height=0.9cm},
  decisione/.style = {diamond, draw, aspect=2, minimum width=2.5cm},
  freccia/.style   = {thick, -{Stealth}}
}

\title{$\mathbb{Z}_p$ e piccolo teorema di Fermat}
\author{Marco Cerbella}
\date{Lezione 14 -- 3 novembre 2026}

\begin{document}
\maketitle

\section{$\mathbb{Z}_p$ con $p$ primo}

\begin{osservazione}
Sia $p$ un numero primo. Allora $[c] \in \mathbb{Z}_p = \{[0], [1], \dots, [p-1]\}$ e, se $[c] \neq [0]$, vale
\[
  \operatorname{MCD}(c, p) = 1 \iff [c] \text{ è invertibile.}
\]
Quindi $(\mathbb{Z}_p, +, \cdot)$ è un \emph{campo}.
\end{osservazione}

\begin{proposizione}
Sia $p$ primo, $p > 0$, e $x, y \in \mathbb{Z}$. Allora
\[
  (x + y)^p \equiv x^p + y^p \pmod p .
\]
\end{proposizione}

\begin{proof}
Per lo sviluppo del binomio di Newton,
\[
  (x + y)^p = \sum_{k=0}^{p} \binom{p}{k} x^k y^{p-k} = x^p + y^p + \sum_{k=1}^{p-1}\binom{p}{k} x^k y^{p-k},
\]
dove si è usato $\binom{p}{p} = 1 = \binom{p}{0}$. Per $1 \le k \le p-1$,
\[
  \binom{p}{k} = \frac{p!}{k!\,(p-k)!} = \frac{p\,(p-1)\cdots(p-k+1)\,(p-k)!}{k!\,(p-k)!} = \frac{p\,(p-1)\cdots(p-k+1)}{k!} .
\]
Il numero $p$ compare al numeratore, ma non divide $k!$ (con $k < p$ e $p$ primo), quindi $p \mid \binom{p}{k}$. Dunque, modulo $p$, tutti i termini della sommatoria sono $\equiv 0$ e
\[
  (x + y)^p \equiv x^p + y^p \pmod p . \qedhere
\]
\end{proof}

\begin{osservazione}
In termini di classi: $\big([x + y]_p\big)^p = \big[(x + y)^p\big]_p = \big[x^p + y^p\big]_p$.
\end{osservazione}

\section{Piccolo teorema di Fermat}

\begin{teorema}[Piccolo teorema di Fermat]
Sia $a \in \mathbb{Z}$ e $p$ primo. Allora
\[
  a^p \equiv a \pmod p .
\]
\end{teorema}

\begin{proof}
Sia dapprima $a \ge 0$: si procede per induzione su $a$ (con $p$ primo, $p > 0$).

\emph{Caso base:} $a = 0$: $0^p \equiv 0 \pmod p$.

\emph{Passo induttivo.} Supponiamo vero $a^p \equiv a \pmod p$ (ipotesi induttiva) e dimostriamo $(a+1)^p \equiv a + 1 \pmod p$. Per la proposizione precedente e per l'ipotesi induttiva,
\[
  (a + 1)^p \equiv a^p + 1^p \equiv a + 1 \pmod p ,
\]
quindi il risultato vale per ogni $a \ge 0$.

Se $a < 0$, allora $(-a) > 0$ e vale il risultato per i numeri non negativi. Si ha
\[
  0 = 0^p = \big(a + (-a)\big)^p \equiv a^p + (-a)^p \equiv a^p + (-a) = a^p - a \pmod p ,
\]
cioè $0 \equiv a^p - a \pmod p$, cioè $a \equiv a^p \pmod p$.
\end{proof}

\begin{teorema}[Corollario]
Se $\operatorname{MCD}(a, p) = 1$ allora $a^{p-1} \equiv 1 \pmod p$.
\end{teorema}

\begin{esempio}
Il resto di $8^3$ diviso $3$: per il piccolo teorema di Fermat $8^3 \equiv 8 \equiv 2 \pmod 3$.

$12\,440^3 \equiv\, ?\pmod 3$: poiché $12\,440 = 3\cdot 4146 + 2$, si ha $12\,440^3 \equiv 12\,440 \equiv 2 \pmod 3$.
\end{esempio}

\begin{esempio}[$4^{194} \bmod 7$]
Usiamo il corollario: $\operatorname{MCD}(4, 7) = 1$, quindi $4^6 \equiv 1 \pmod 7$. Dividiamo $194$ per $6$: $194 = 6\cdot 32 + 2$. Allora
\[
  4^{194} = 4^{6\cdot 32 + 2} = \left(4^6\right)^{32}\cdot 4^2 \equiv 1^{32}\cdot 4^2 = 16 \equiv 2 \pmod 7 .
\]
\end{esempio}

\begin{osservazione}
Il corollario ha l'ipotesi $\operatorname{MCD}(a, p) = 1$. Per $a = 6$, $p = 3$: $6^3 \equiv 6 \pmod 3$ (piccolo teorema di Fermat) è vero, ma poiché $\operatorname{MCD}(6, 3) \neq 1$, la congruenza $6^2 \equiv 1 \pmod 3$ è \emph{falsa}: $6^2 = 36 \equiv 0 \pmod 3$.
\end{osservazione}

\end{document}
